\documentclass{article}
\usepackage[T1]{fontenc}
\usepackage{lmodern}
\usepackage{graphicx}
\usepackage{amsmath,amssymb}
\usepackage[a4paper,margin=2cm]{geometry}
\usepackage[absolute,overlay]{textpos}
\setlength{\TPHorizModule}{1mm}
\setlength{\TPVertModule}{1mm}
\usepackage[indonesian]{babel}
\usepackage{hyperref}
\setlength{\emergencystretch}{2em}
% unit: Halaman 44

\begin{document}

% === Halaman 44 (llm) ===

\section*{Elliptic Curve Elgamal (ECEG)}

\begin{itemize}
    \item \textit{Elliptic Curve Elgamal}: sistem kriptografi kurva eliptik yang diadopsi dari algoritma El Gamal.
    \item Misalkan \textcolor{red}{Alice} ingin mengirim \textcolor{red}{Bob} pesan M yang dienkripsi.
    \begin{itemize}
        \item Baik Alice dan Bob menyepakati kurva elelitik dan titik basis B.
        \item Alice dan Bob membuat kunci privat/kunci publik.
        \begin{itemize}
            \item Alice
            \begin{itemize}
                \item Kunci privat = a
                \item Kunci publik = $P_A = a \cdot B$
            \end{itemize}
            \item Bob
            \begin{itemize}
                \item Kunci privat = b
                \item Kunci publik = $P_B = b \cdot B$
            \end{itemize}
        \end{itemize}
    \end{itemize}
    \item Alice mengambil plainteks, M, lalu mengkodekannya menjadi sebuah titik, $P_M$, pada kurva eliptik
\end{itemize}

\vspace{5mm}

\textcolor{red}{*) Sumber bahan: Debdeep Mukhopadhyay, Elliptic Curve Cryptography , Dept of Computer Sc and Engg IIT Madras}

\end{document}
